\section{Agente de IA: planificación, razonamiento y uso de herramientas}

\subsection{Qué es un Agente de IA}

El docente señala que el término «AI Agent» tiene significados muy distintos según quién lo use: desde una simple llamada a la API de un LLM hasta sistemas complejos de colaboración entre varios LLM se llaman «Agent». Según el docente, las dos dimensiones de capacidad más centrales de un Agente son:

\begin{enumerate}
    \item \textbf{Planificación y razonamiento} (Planning \& Reasoning): capacidad de descomponer tareas complejas, elaborar un plan y reflexionar sobre resultados intermedios
    \item \textbf{Uso de herramientas} (Tool Use): capacidad de llamar a APIs externas (motores de búsqueda, calculadora, bases de datos, etcétera) para obtener información o ejecutar acciones
\end{enumerate}

\subsection{ReAct: la fusión de razonamiento y acción}

El docente presentó el artículo de ReAct (Reasoning + Acting), uno de los trabajos fundacionales más influyentes en el área de los Agentes.

Antes de ReAct, quienes investigaban hacían que el modelo solo razonara (reasoning traces) o solo ejecutara acciones (action traces). ReAct combina ambas cosas en un marco híbrido.

\textbf{Comparación concreta}:

\textit{Tarea}: determinar si es verdadero que ``Reign Over Me is an American film made in 2010''.

\begin{table}[H]
\centering
\begin{tabular}{p{6cm}p{6cm}}
\toprule
\textbf{Chain-of-Thought (solo razonamiento)} & \textbf{ReAct (razonamiento + acción)} \\
\midrule
Thought: It is an American film... & Thought: I need to search for Reign Over Me \\
Thought: It was made in 2010... & Action: Search{[}Reign Over Me{]} \\
Answer: \textbf{True} {\color{red}(incorrecto)} & Observation: American film, released \textbf{2007} \\
 & Thought: It was made in 2007, not 2010 \\
 & Answer: \textbf{False} {\color{green!60!black}(correcto)} \\
\bottomrule
\end{tabular}
\caption{Comparación entre Chain-of-Thought y ReAct}
\end{table}

El método CoT puro, al razonar solo con base en el «conocimiento interno» del modelo, concluyó erróneamente que la película se filmó en 2010 (en realidad fue en 2007). ReAct obtuvo el año de estreno correcto al buscar en una fuente de información externa y así llegó a la respuesta correcta.

\begin{importantbox}{La estructura cíclica del marco ReAct}
El núcleo de ReAct es el ciclo Thought $\rightarrow$ Action $\rightarrow$ Observation:
\begin{enumerate}
    \item \textbf{Thought}: el modelo razona sobre el estado actual y decide qué hacer a continuación
    \item \textbf{Action}: el modelo ejecuta una acción (como buscar o llamar a una API)
    \item \textbf{Observation}: el entorno devuelve el resultado de la acción
    \item Se repite el ciclo anterior hasta que el modelo considera que tiene información suficiente para dar la respuesta final
\end{enumerate}
Este marco se ha convertido en la arquitectura base de los sistemas de Agentes modernos (como LangChain, AutoGPT, etcétera).
\end{importantbox}

El docente mostró también otro ejemplo de ReAct: un juego de aventura de texto cuya tarea es «poner el frasco de pimienta sobre el cajón».

\begin{itemize}
    \item Modo \textbf{solo acción}: el Agente camina hasta el fregadero 1, intenta tomar un frasco de pimienta que no existe ahí y queda atrapado en un ciclo
    \item Modo \textbf{ReAct}: el Agente razona que «el fregadero 1 no tiene frasco de pimienta, hay que buscar en otro lugar», navega con éxito y completa la tarea
\end{itemize}

\subsection{Toolformer: enseñar a un LLM a usar herramientas}

El segundo artículo fundacional que presentó el docente es Toolformer, de Meta, que enseña a un LLM a llamar de manera autónoma a APIs externas mientras genera texto.

Los tipos de herramientas que admite Toolformer incluyen:
\begin{itemize}
    \item \textbf{Herramienta de QA}: responde preguntas factuales
    \item \textbf{Calculator}: ejecuta cálculos precisos
    \item \textbf{Machine Translation}: traduce
    \item \textbf{Wikipedia Search}: busca conocimiento enciclopédico
\end{itemize}

\begin{knowledgebox}{El método de entrenamiento de Toolformer: un pipeline de tres pasos}
\begin{enumerate}
    \item \textbf{Generación de anotaciones}: se le dan al modelo unos pocos ejemplos (few-shot) para que anote en el texto «aquí debería llamarse a cierta API»
    \item \textbf{Ejecución de la API}: se llama de verdad a esas APIs y se obtienen los resultados que devuelven
    \item \textbf{Filtrado de eficacia} (este es el paso más ingenioso): se compara la loss de entrenamiento entre el caso «con el resultado de la API» y el caso «sin el resultado de la API». Solo se conservan como datos de entrenamiento las llamadas a la API que \textbf{redujeron notablemente la loss}
\end{enumerate}
Por ejemplo, la llamada a la API en ``Pittsburgh was known as the {[}QA: Steel City{]}'' resulta útil (el modelo originalmente no estaba muy seguro del apodo de Pittsburgh), mientras que la llamada en ``the country that Pittsburgh is in is {[}QA: United States{]}'' no aporta nada (el modelo ya lo sabía de antemano), por lo que esta última se filtra y se descarta.
\end{knowledgebox}

\subsection{Resumen de la sección}

Las dos capacidades centrales de un Agente de IA —la planificación y el razonamiento, y el uso de herramientas— fueron desarrolladas respectivamente por dos artículos fundacionales: ReAct y Toolformer. ReAct fusiona el razonamiento (Thought) y la acción (Action) en un marco cíclico que permite al Agente ajustar su estrategia de manera dinámica; Toolformer, por su parte, le enseña al modelo cuándo y cómo llamar a herramientas externas, y garantiza mediante un mecanismo de filtrado ingenioso que solo se aprendan patrones de llamada a herramientas que aporten valor.
